\DocumentMetadata{
  pdfversion=2.0,
  pdfstandard=ua-2,
  lang=en-US,
  tagging=on,
  tagging-setup={math/setup=mathml-SE,tabsorder=structure}
}
\documentclass[11pt,letterpaper]{article}
\usepackage[margin=0.68in,includefoot]{geometry}
\usepackage{newcomputermodern}
\usepackage{amsmath,amscd,mathtools}
\usepackage{tikz,adjustbox}
\providecommand{\square}{\mdlgwhtsquare}
\usepackage[hidelinks]{hyperref}
\hypersetup{
  pdftitle={Analysis First-Year Exam, May 2022, Part 2},
  pdfauthor={University of Florida Department of Mathematics},
  pdfsubject={Graduate mathematics examination},
  pdfdisplaydoctitle=true,
  bookmarksopen=true
}
\setcounter{secnumdepth}{0}
\setlength{\parindent}{0pt}
\setlength{\parskip}{0.28em}
\setlength{\emergencystretch}{3em}
\setlength{\leftmargini}{2.15em}
\setlength{\leftmarginii}{2.65em}
\setlength{\leftmarginiii}{3em}
\renewcommand{\labelenumi}{\textbf{\theenumi.}}
\pagestyle{empty}
\raggedbottom
\allowdisplaybreaks
\newenvironment{examproblems}[1]{%
  \begin{enumerate}%
  \setcounter{enumi}{#1}%
  \setlength{\topsep}{0.35em}%
  \setlength{\partopsep}{0pt}%
  \setlength{\itemsep}{0.48em}%
  \setlength{\parsep}{0.18em}%
}{\end{enumerate}}
\newenvironment{examsubparts}{%
  \begin{enumerate}%
  \setlength{\topsep}{0.18em}%
  \setlength{\partopsep}{0pt}%
  \setlength{\itemsep}{0.16em}%
  \setlength{\parsep}{0.1em}%
}{\end{enumerate}}
\newenvironment{examinstructions}{%
  \par\begingroup\itshape
}{%
  \par\endgroup\vspace{0.18em}
}
\makeatletter
\renewcommand\section{\@startsection{section}{1}{0pt}%
  {-1.25ex plus -.25ex minus -.1ex}%
  {0.55ex plus .1ex}%
  {\normalfont\large\bfseries}}
\renewcommand{\@maketitle}{%
  \newpage\null\vskip -1em%
  {\footnotesize\bfseries
    \href{https://www.ufl.edu/}{University of Florida}, \href{https://math.ufl.edu/}{Department of Mathematics}\par}%
  \vskip -0.5em%
  \tagstructbegin{tag=Title}%
    {\LARGE\bfseries\href{https://gma.math.ufl.edu/exams/analysis-first-year/}{Analysis First-Year Exam}, May 2022, Part 2\par}%
  \tagstructend%
  \vskip 0.25em%
  \tagmcbegin{artifact}%
    \hrule height 1pt%
  \tagmcend%
  \vskip 1em%
}
\makeatother
\title{Analysis First-Year Exam, May 2022, Part 2}
\author{}
\date{}
\begin{document}
\maketitle
\thispagestyle{empty}
\begin{examinstructions}
Answer FOUR questions in detail. State carefully any results used without proof.
\end{examinstructions}
\begin{examproblems}{0}
\item
Let \(f:[a,b]\to\mathbb{R}\) have the property that for each \(\varepsilon>0\) there exist Riemann integrable functions~\(\ell\) and~\(u\) on \([a,b]\) such that \(\ell\leq f\leq u\) and \(\int_a^b(u-\ell)<\varepsilon\). Does it follow that~\(f\) is Riemann integrable on \([a,b]\)? Proof or counterexample.
\item
Suppose \((p_n)_{n=0}^{\infty}\) is a sequence of polynomials in one variable.
\begin{examsubparts}
\item[\textnormal{(i)}]
Assume that \(p_n\to f\) uniformly on \([0,1]\) as \(n\to\infty\); deduce as much as possible about the function \(f:[0,1]\to\mathbb{R}\).
\item[\textnormal{(ii)}]
Assume that \(p_n\to f\) uniformly on~\(\mathbb{R}\) as \(n\to\infty\); deduce as much as possible about the function \(f:\mathbb{R}\to\mathbb{R}\).
\end{examsubparts}
\item
Let the continuous function \(f:[0,1]\to\mathbb{R}\) satisfy 
\[
\int_0^1 f(t)t^n\,dt=1/(n+2)
\]
 for all but finitely many positive integers~\(n\). Deduce as much as possible about~\(f\).
\item
Let \((f_n)_{n=0}^{\infty}\) be a sequence of measurable real-valued functions on some measurable space. Prove that each of the following sets is measurable:
\begin{examsubparts}
\item[\textnormal{(i)}]
\(A=\{\omega:\sum_{n=0}^{\infty}f_n(\omega)\text{ is absolutely convergent}\}\);
\item[\textnormal{(ii)}]
\(C=\{\omega:\sum_{n=0}^{\infty}f_n(\omega)\text{ is conditionally convergent}\}\).
\end{examsubparts}
\item
Let \((\Omega,\mathcal{A},\mu)\) be a measure space; let \(A_0\subseteq A_1\subseteq\cdots\) be an increasing sequence in~\(\mathcal{A}\) such that \(\bigcup_{n=0}^{\infty}A_n=\Omega\); let \(f:\Omega\to\mathbb{R}\) be measurable; and let~\(L\) be a real number. Consider the statements:

\par
Prove that (i) implies (ii) and decide (with proof or counterexample) whether (ii) implies (i).
\begin{examsubparts}
\item[\textnormal{(i)}]
\(f\) is integrable on~\(\Omega\) and \(\int_\Omega f\,d\mu=L\);
\item[\textnormal{(ii)}]
\(f\) is integrable on each \(A_n\) and \(\int_{A_n}f\,d\mu\to L\) as \(n\to\infty\).
\end{examsubparts}
\end{examproblems}
\end{document}
